\documentclass{article}

% Language setting
% Replace `english' with e.g. `spanish' to change the document language
\usepackage[english]{babel}

% Set page size and margins
% Replace `letterpaper' with `a4paper' for UK/EU standard size
\usepackage[letterpaper,top=2cm,bottom=2cm,left=3cm,right=3cm,marginparwidth=1.75cm]{geometry}

% Useful packages
\usepackage{amsmath}
\usepackage{graphicx}
\usepackage[colorlinks=true, allcolors=blue]{hyperref}

\title{Labwork 1: GradientDescent}
\author{Nguyen Nhu Duy - 2540044}

\begin{document}
\setlength{\parindent}{0pt}
\maketitle
\section{Implement the algorithm}
Step 1: define f(x) and f\_(x)

Step 2: Random initialization: $x = x_0$, define the learning rate r, and the threshold to stop 

Step 3: define the iteration follow the below pseudo code

\begin{itemize}
    \item  Assign:  $x = x - r \times f\_(x)$ (where r is learning rate)
    \item  Re-compute f(x). Stop if f(x) is small enough, or repeat step 2 if not 1
\end{itemize}

\section{Effect of different learning rate r}
\begin{itemize}
    \item Small $r$: Slow convergence; requires excessive iterations.
    \item Large $r$: Risk of overshooting; fails to reach the minimum.
    \item Optimal $r$: Ensures swift and stable convergence to the minimum.
\end{itemize}

\end{document}